\section{Introduction}
\label{sec:introduction}

For three decades, a new vehicle function usually meant a new electronic
control unit (ECU). Premium cars already carried 70 to 100 ECUs in 2009
\cite{charette2009car}, with a wiring harness, a supply chain and an
integration effort to match \cite{bandur2021centralized,
askaripoor2022synthesis}. The industry is now consolidating this software
onto a few powerful computers. A \emph{central computer} runs most
functions, and a handful of \emph{zone controllers} connect the sensors and
actuators of their physical area of the car. An Ethernet backbone with
Time-Sensitive Networking (TSN) links them \cite{lobello2019perspective,
ieee8021q2022}. This \emph{zonal} electrical/electronic (E/E) architecture
shortens the harness and decouples software from hardware
\cite{bandur2021centralized,mauser2025centralization}.

It also turns a question that each ECU supplier used to answer alone into
a system-wide one: the design question in the box.
\begin{figure}[t]
\begin{keybox}[unbreakable,before skip=0pt,after skip=0pt]{Design question}
How many processor cores does the software need, and on which ECU? Does
every function still meet its timing requirement once its tasks are spread
over cores and over the network? And how much room is left for the
functions that software updates will add during the life of the vehicle?
\end{keybox}
\end{figure}

\paragraph{Why the question is hard.} Vehicle functions are
\emph{task-chains} (also called cause-effect chains): a sensor is sampled,
data flows through a sequence of tasks, and an actuator is commanded
\cite{feiertag2008compositional,kramer2015waters}. Their timing
requirement is end to end, typically a bound on the \emph{maximum reaction
time}: the longest time from an external event to the actuation that
accounts for it \cite{feiertag2008compositional,durr2019end,
gunzel2023equivalence}.
On a zonal platform this reaction time depends on four decisions, and they
interact.
\begin{enumerate}[label=(\roman*)]
\item The \emph{allocation} of tasks to cores. It is constrained by the
wiring (a sensor driver runs in the zone of its sensor), by memory, by
safety (redundant functions on different ECUs, safety tasks on lock-step
cores) and by cores of different types and speeds.
\item The \emph{priorities} on each core, which set the response times of
the tasks.
\item The \emph{communication between cores}. With the Logical Execution
Time (LET) paradigm \cite{kirsch2012let,hamann2017communication,
biondi2018achieving}, data crosses a core boundary only at fixed instants.
This makes timing predictable, but every crossing adds sampling delays and
every core pays for the copy task.
\item The \emph{network}. Every task-chain that crosses ECUs adds a stream
of frames, which shares output ports with the other streams. Its delay
depends on the traffic shaping of the switches \cite{ieee8021q2022,
zhao2022quantitative}.
\end{enumerate}
Packing tasks onto fewer cores saves hardware but lengthens response
times. Spreading a task-chain over ECUs frees cores but adds LET sampling
delays and network streams, and these streams delay other task-chains on
the same ports. The platform also has to survive its own life: the
hardware is frozen at the start of production (SOP), while software
updates keep adding functions for years.

\paragraph{What exists.} Each piece is well studied on its own.
Allocation and priority assignment for distributed real-time systems goes
back to Tindell et al.~\cite{tindell1992allocating} (see
\cite{zhu2012optimization,davis2011survey,davis2016review}). The
end-to-end analysis of task-chains has matured \cite{becker2017end,
durr2019end,gunzel2021timing,martinez2018analytical}. LET deployment on
multicores has been optimized \cite{biondi2018achieving,
pazzaglia2019optimizing,paladino2024optimizing}. TSN has well-established
Network Calculus bounds \cite{azua2014complete,mohammadpour2018latency,
zhao2021latency,zhao2018worst,leboudec2018theory}, multicore
response-time analysis accounts for shared-memory interference
\cite{davis2018extensible}, and compositional system-level analysis
chains the analyses of processors and buses for a given mapping
\cite{henia2005system}. Studies of zonal
architectures, however, focus on the harness and the network
\cite{park2024zonal} or schedule tasks without a schedulability analysis
\cite{aghababaiyan2025deterministic}, and syntheses of E/E architectures
survey the mapping of software without a system-level timing analysis
\cite{askaripoor2022synthesis}.
What is missing is the system level. To our knowledge, and as
\cref{sec:related} details, no published method decides the allocation of
an automotive application while it analyses, together and with the same
granularity, the fixed-priority scheduling of every core, the
shared-memory contention between its cores, LET-based communication and the
delays of shaped TSN streams, under the spatial constraints of a vehicle,
and sizes the cores from that analysis.

\paragraph{Contributions.} This paper gives such a system-level method,
in enough detail to re-implement it from the text, with the code as a
reference.
\begin{enumerate}[label=\textbf{C\arabic*.}, leftmargin=2.2em]
\item \textbf{A formal model and problem statement}
(\cref{sec:model,sec:problem}). The model covers zonal, domain, hybrid
and distributed E/E architectures (\cref{sec:other-architectures}),
heterogeneous cores,
a utilization cap, communication between cores at LET instants (implicit
or strict LET), multi-rate task-chains,
shared-memory contention, and a TSN backbone with
four TSN mechanisms: credit-based shaping, time-aware shaping,
asynchronous traffic shaping and frame replication, on a star, a tree, a
ring or any switched topology. It also covers the
spatial constraints of a vehicle: allowed ECUs, memory, separation of
redundant task-chains, co-location and allowed core types. The objective
is to use the fewest cores, then to meet as many end-to-end deadlines as
possible, optionally for a target timing margin.
\item \textbf{Timing analyses with proofs} (\cref{sec:analysis}). They are
an optimal priority assignment that respects precedence, response times
with shared-memory contention, a bound on the maximum reaction time of a
task-chain under implicit communication and under strict LET, Network
Calculus bounds for the four TSN mechanisms,
and a lower bound on the number of cores of any mapping of the
sub-chains.
\item \textbf{A mapping heuristic} (\cref{sec:method}), given as
pseudo-code. Its new parts are a communication-aware allocation rule, a
core consolidation that preserves schedulability, a reservation of zone
memory for pinned tasks, a deadline-driven local search with an exact
incremental re-analysis, and the sizing of the platform for a target
margin.
\item \textbf{An exact formulation} (\cref{sec:exact}): a constraint
program that finds the fewest cores, with schedulability added as lazy
constraints. It measures the optimality gap of the heuristic for the
sub-chains of the partitioning, without contention.
\item \textbf{A lifecycle analysis} (\cref{sec:lifecycle}). The platform
is sized at the SOP, and then software updates add task-chains to the
fixed hardware. This trades the hardware cost against the number of
updates the platform absorbs.
\item \textbf{An evaluation} (\cref{sec:evaluation}). The synthetic
applications follow the statistics of an industrial engine-management
software \cite{kramer2015waters}. The evaluation adds a small detection
chain and a whole-vehicle application of 2000 task-chains. A frame-level network simulation and a full-system simulation
check the bounds. The method is compared with the optimum, with simulated
annealing and a genetic algorithm, and with the rules of practice; each of
its parts is removed in turn; strict LET is compared with implicit
communication; the same software is mapped onto four architectures; and
each assumption of the generator is varied.
\end{enumerate}
\begin{keybox}{Central claim}
A system-level method that couples allocation, fixed-priority scheduling,
shared-memory contention, LET and TSN bounds sizes the cores of a vehicle
platform close to the optimal allocation of its sub-chains (2.3 to
2.9\,\% above it on average, proven optimal in 81\,\% of 140 instances)
in milliseconds, and no end-to-end bound it reported was exceeded when
the network and the whole system were simulated.
\end{keybox}
The method is implemented as an open, deterministic C library and
command-line tool, which reads the application and the platform from JSON
files. A separate user manual describes the tool, five use cases, and how
to reproduce every number of this paper.

\paragraph{How to read this paper.} \Cref{sec:background} recalls the
terms and mechanisms the paper builds on. A reader who knows LET, TSN and
Network Calculus can skip it. \Cref{sec:related} reviews the state of the
art. \Cref{sec:model,sec:problem} define the model and the problem.
\Cref{sec:analysis} gives the timing analyses, \cref{sec:method} the
heuristic, \cref{sec:exact} the exact formulation and
\cref{sec:lifecycle} the lifecycle analysis. \Cref{sec:example} follows
one task-chain through all the steps. \Cref{sec:evaluation} reports the
experiments, \cref{sec:discussion} discusses design rules and limits, and
\cref{sec:conclusion} concludes. Every section starts with a short summary
of what it contains.
